% Bagean: metodhe
% Bab: bukti
% Pérangan: conto-2

\documentclass[../../../include/open-logic-section]{subfiles}

\begin{document}

\olfileid{mth}{prf}{ex2}

\olsection{Conto Liyané}

\begin{prop}
Yèn $A \subseteq C$, mula $A \cup (C \setminus A) = C$.
\end{prop}

\begin{proof}
  Anggepen $A \subseteq C$. Kita arep mbuktekaké
  $A \cup (C \setminus A) = C$.
  \begin{quote}
    Wiwitané, pratelan iki awujud kondisional. Kanthi
    implisit, pratelan iki dikuantifikasi universal: proposisi
    mau bener kanggo kabèh himpunan $A$ lan $C$. Dadi $A$ lan
    $C$ iku variabel kanggo himpunan sembarang. Kanggo
    mbuktekaké pratelan kaya iki, kita nganggep antesedèn
    lan mbuktekaké konsekuèné.

    Sabanjuré gunakna anggepan $A \subseteq C$.
    Dhéfinisi~$\subseteq$ dijlentrehaké: saben !!{element}s
    saka~$A$ uga !!{element}s saka~$C$. Kasunyatan wigati
    iki ayo ditulis, amarga bakal digunakaké ing sadawaning
    bukti.
  \end{quote}
  Miturut dhéfinisi~$\subseteq$, amarga $A \subseteq C$,
  kanggo saben $z$, yèn $z \in A$, mula $z \in C$.
  \begin{quote}
    Kabèh dhéfinisi ing anggepan wis dijlentrehaké. Saiki
    pindhah menyang kesimpulan. Kita arep nuduhaké
    $A \cup (C \setminus A) = C$, mula bukti disusun kaya
    conto sadurungé: saben !!{element} saka
    $A \cup (C \setminus A)$ uga !!a{element} saka~$C$,
    lan kosok baliné saben !!{element} saka $C$ uga
    !!a{element} saka $A \cup (C
    \setminus A)$. Iki bisa dicekak dadi:
    $A \cup (C \setminus A)
    \subseteq C$ lan $C \subseteq A \cup (C \setminus A)$.
    Ing kéné kita tumindak kosok baliné saka njlentrehaké
    dhéfinisi, supaya buktiné luwih gampang diwaca.
    Amarga iki konjungsi, loro pérangan kudu dibuktèkaké.
    Kanggo pérangan kapisan, yaiku yèn saben
    !!{element} saka $A \cup (C \setminus A)$ uga
    !!a{element} saka~$C$, kita nganggep
    $z \in A \cup (C \setminus A)$ kanggo~$z$ sembarang
    lan nuduhaké $z \in C$. Miturut dhéfinisi
    $\cup$, kita nyimpulaké $z \in A$ utawa
    $z \in C \setminus A$ saka
    $z \in A \cup (C \setminus A)$.
    Saiki panjenengan mesthiné wis wiwit kulina.
  \end{quote}
  $A \cup (C \setminus A) = C$ yèn lan mung yèn
  $A \cup (C \setminus A) \subseteq
  C$ lan $C \subseteq A \cup (C \setminus A)$.
  Kaping pisan kita buktèkaké
  $A \cup (C \setminus A) \subseteq C$.
  Jupuken $z \in A \cup (C
  \setminus A)$. Dadi $z \in A$ utawa
  $z \in (C \setminus A)$.
  \begin{quote}
    Kita tekan disjungsi, lan saka kéné arep
    mbuktekaké $z \in C$. Gunakna bukti mawa kasus.
  \end{quote}
  Kasus 1: $z \in A$. Amarga kanggo saben $z$,
  yèn $z \in A$ mula $z \in C$, kita olèh $z \in C$.
  \begin{quote}
    Ing kéné kita migunakaké kasunyatan kang wis
    dicathet saka hipotesis proposisi $A \subseteq C$.
    Kasus kapisan rampung; saiki kasus kapindho,
    $z \in (C
    \setminus A)$. Élinga yèn $C \setminus A$ iku
    \emph{selisih} rong himpunan: himpunan kabèh
    !!{element}s saka~$C$ kang dudu !!{element}s saka~$A$.
    Saben !!{element} saka $C$ kang ora ana ing~$A$
    mesthi uga !!a{element} saka~$C$.
  \end{quote}
  Kasus 2: $z \in (C \setminus A)$. Iki ateges
  $z \in C$ lan $z
  \notin A$. Mula, mliginé, $z \in C$.
  \begin{quote}
    Apik, arah kapisan wis kabuktèkaké. Saiki
    arah kapindho: buktèkna
    $C \subseteq A \cup (C \setminus
    A)$. Mula kita nganggep $z \in C$ lan mbuktekaké
    $z \in A \cup (C
    \setminus A)$.
  \end{quote}
  Saiki jupuken $z \in C$. Kita arep nuduhaké
  $z \in A$ utawa $z \in C
  \setminus A$.
  \begin{quote}
    Amarga saben !!{element}s saka $A$ uga
    !!{element}s saka $C$, lan $C
    \setminus A$ yaiku himpunan kabèh obyèk kang
    !!{element}s saka $C$ nanging dudu anggota $A$,
    mula $z$ mesthi ana ing $A$ utawa $C
    \setminus A$. Iki mbokmenawa durung cetha
    yèn panjenengan durung ngerti sebabé. Luwih becik
    mbuktekaké langkah demi langkah. Kita bisa
    migunakaké kasunyatan prasaja tanpa bukti:
    $z \in A$ utawa $z \notin A$. Iki diarani
    ``Hukum Tengah kang Disingkiraké'': kanggo
    sembarang pratelan~$p$, $p$ bener utawa
    negasiné bener. Ing kéné, $p$ yaiku pratelan
    $z \in A$. Amarga iki disjungsi, kita bisa
    migunakaké bukti mawa kasus maneh.
  \end{quote}
  Salah siji $z \in A$ utawa $z \notin A$.
  Ing kasus kapisan, $z \in A \cup
  (C \setminus A)$. Ing kasus kapindho,
  $z \in C$ lan $z \notin A$, mula
  $z \in C \setminus A$. Dadi
  $z \in A \cup (C \setminus A)$.
  \begin{quote}
    Bukti wis rampung: kita wis nuduhaké
    $A \cup (C \setminus A) = C$.
  \end{quote}
\end{proof}

\end{document}
